\begin{table}
\caption{Magnetic-ordering performance with and without compounds lacking a magnetism-capable element}
\label{tab:tableS7_ordering_decomposition}
\tablefont
\begin{tabularx}{\textwidth}{@{}>{\raggedright\arraybackslash}Xr@{}}
\toprule
\textbf{Quantity} & \textbf{Mean over 10 splits} \\
\midrule
Macro-F1 & 0.692 \\
Macro-F1 (compounds with a magnetism-capable element) & 0.681 \\
Share of test rows with a magnetism-capable element & 0.683 \\
Share of non-magnetic rows with a magnetism-capable element & 0.388 \\
Recall, FM & 0.774 \\
Recall, FM (compounds with a magnetism-capable element) & 0.800 \\
Recall, AFM & 0.546 \\
Recall, AFM (compounds with a magnetism-capable element) & 0.551 \\
Recall, FiM & 0.485 \\
Recall, FiM (compounds with a magnetism-capable element) & 0.500 \\
Recall, NM & 0.906 \\
Recall, NM (compounds with a magnetism-capable element) & 0.837 \\
Macro-F1 over FM, AFM and FiM (compounds with a magnetism-capable element) & 0.654 \\
\bottomrule
\end{tabularx}
\par\vspace{4pt}\raggedright\noindent Tier 1 XGBoost on the full ordering population, 10 prespecified splits; magnetism-capable element = curation rule C1 list.
\end{table}
